\subsection{Interpretation of Results}

Several important patterns across different compression paradigms are revealed by the experimental results obtained in **Part 1** of this ongoing project.

\paragraph{Entropy Coding (Huffman, Arithmetic, LZW)}  
The near‑entropy performance of Huffman (1.07:1) and arithmetic coding (1.08:1) confirms that first‑order entropy models alone are insufficient for natural images. The original image entropy (7.44 bpp) is reduced only marginally. LZW (1.15:1) outperformed both by capturing repeating pixel patterns (e.g., in the textured hat of Lenna), which first‑order statistics ignore. However, the improvement is modest, indicating that spatial redundancy is the dominant factor.

\paragraph{Spatial Prediction (JPEG Lossless, LOCO‑I, CALIC)}  
The significant entropy reduction achieved by LOCO‑I (5.60 bpp, 24.8% savings) demonstrates the power of context‑based prediction. The median‑edge detector (MED) adapts to local gradients, producing residuals heavily peaked at zero. JPEG Lossless (mode 4, 5.99 bpp) performed worse because its planar predictor \(a+b-c\) fails at edges. CALIC (5.76 bpp) lies between the two; its full context modelling would likely excel on highly textured images. These results validate that adaptive predictors are essential for lossless compression.

\paragraph{Transform Domain (S‑Transform, S+P, Lifting)}  
All three integer‑to‑integer transforms achieved perfect losslessness with zero reconstruction error. However, entropy reductions were modest (≈0.56 bpp), suggesting that simple Haar‑like wavelets do not compact energy as effectively as more sophisticated transforms (e.g., Daubechies) that are not integer‑to‑integer. The S+P transform did not outperform the basic S‑Transform in our implementation, likely due to the fixed prediction coefficients. This highlights the sensitivity of predictive wavelet transforms to coefficient tuning.

\paragraph{Lossy+Residual Hybrid}  
The hybrid scheme achieved a lossless compression ratio of 1.54:1, outperforming the best pure lossless method (LZW) by 40\%. The key insight is that the residual after a good lossy approximation has very low entropy (3.58 bpp). This validates two‑stage lossy+residual coding as a highly effective strategy for lossless compression when the lossy stage is carefully chosen.

\paragraph{Colour Quantisation (K‑Means, PSO, Gath‑Geva)}  
K‑Means produced high compression ratios (up to 23.98:1 for Cameraman with K=2, and 5.99:1 for Lenna with K=16). As expected, the objective function \(J(V)\) decreased with increasing \(K\). PSO achieved a lower MSE than K‑Means for the same \(K\) (e.g., 126.8 vs. unreported) but at a much higher computational cost (24.8 s vs. 0.64 s), reflecting the classic exploration–exploitation trade‑off. Gath‑Geva with \(K=8\) gave a compression ratio of 7.99:1 on Lenna, demonstrating that ellipsoidal clusters handle colour distributions more flexibly than spherical K‑Means.

\paragraph{Fuzzy Backpropagation Neural Network (FBPNN)}  
The network achieved a compression ratio of 5.97:1 with 15 colours. Training loss dropped smoothly from 0.07734 to 0.05505 over 30 epochs, indicating stable convergence. The fuzzy weighting of error updates likely helped avoid local minima. However, the final MSE was only slightly better than standard K‑Means with the same \(K\), suggesting that the added complexity of fuzzy neural learning may not be fully justified for simple colour quantisation tasks – ambiguous data would be required to realise its full potential.

\subsection{Strengths of the Approach (Part 1)}

\begin{itemize}
    \item \textbf{Comprehensive comparative framework:} A wide range of algorithms – from classical entropy coders to modern neural and metaheuristic methods – were implemented under a unified evaluation protocol, enabling direct, fair comparisons.
    \item \textbf{Modular, transparent implementations:} Every algorithm was coded from scratch, avoiding high‑level ML libraries. This ensures full control over parameters, eliminates hidden biases, and makes results highly reproducible and explainable.
    \item \textbf{Adaptive and context‑aware techniques:} LOCO‑I and CALIC demonstrated that exploiting local image structure (gradients, texture patterns) yields substantial entropy reductions (up to 25\% compared to first‑order entropy coding).
    \item \textbf{Hybrid lossless/lossy strategies:} The lossy+residual method achieved a lossless ratio (1.54:1) unattainable by pure lossless methods, showing that carefully designed hybrid schemes can bridge the gap between lossy and lossless coding.
    \item \textbf{Optimisation flexibility:} PSO and Gath‑Geva provide tunable trade‑offs between compression quality and speed. PSO can be run with many iterations for near‑optimal palettes (offline archival), while K‑Means or Gath‑Geva with few clusters can be chosen for real‑time use.
\end{itemize}

\subsection{Limitations (Addressed in Part 2)}

Several limitations identified in Part 1 will be explicitly addressed in the continuation of this project (Part 2, next semester):

\begin{itemize}
    \item \textbf{Scalability of PSO and FBPNN:} Training time for PSO grows linearly with the number of clusters and particles (≈40 s for \(K=64\) on a 256×256 image), rendering real‑time applications impractical. Similarly, FBPNN required several seconds even for a 128×128 image. Part 2 will explore GPU acceleration (JAX, Metal) and adaptive swarm sizing to reduce runtime.
    \item \textbf{Suboptimal transform coefficients:} The fixed prediction coefficients used in the S+P transform (e.g., \(a=[0.5,0.25,0.125]\)) are not optimised for specific images. Future work will investigate learning‑based coefficient selection to improve energy compaction.
    \item \textbf{Limited texture handling in LOCO‑I:} Although LOCO‑I works well on smooth and edge regions, it still produces larger residuals in highly textured areas (e.g., Lenna’s hat). Part 2 will integrate more sophisticated context models (inspired by CALIC) while keeping complexity manageable.
    \item \textbf{Colour space limitation:} All colour quantisation methods operated in RGB space, which is not perceptually uniform. Part 2 will experiment with perceptually uniform spaces (e.g., CIELAB) and incorporate perceptual metrics (SSIM, MS‑SSIM) into the fitness functions.
    \item \textbf{Single image focus:} The current experiments focused mainly on Lenna and Cameraman. Part 2 will evaluate the algorithms on a larger, more diverse dataset including medical, satellite, and synthetic images to improve generalisability.
\end{itemize}

\subsection{Threats to Validity}

\paragraph{Internal Validity (Experimental Design)}
\begin{itemize}
    \item \textbf{Implementation correctness:} All algorithms were verified against known outputs (e.g., zero reconstruction error for lossless methods). Minor implementation differences from reference standards (e.g., the exact distance formula in Gath‑Geva) could affect comparative performance. Core equations were taken directly from original papers to mitigate this.
    \item \textbf{Parameter selection:} Parameters (e.g., \(c_1, c_2, w\) for PSO; \(m\) for fuzzy clustering) were chosen based on typical literature values, not systematically optimised per image. Some algorithms may have been disadvantaged. A full grid search would improve fairness but is computationally expensive.
\end{itemize}

\paragraph{External Validity (Generalizability)}
\begin{itemize}
    \item \textbf{Test images:} Only two natural images (Lenna, Cameraman) were used. Other image classes (medical, synthetic, high‑noise) may not be well represented. Part 2 will include a larger, more diverse dataset.
    \item \textbf{Image size:} Most experiments used 128×128 or 256×256 images for speed. Relative performance of \(O(N \log N)\) vs. \(O(NK)\) algorithms might change with larger images (e.g., 1024×1024).
\end{itemize}

\paragraph{Construct Validity (Measurement)}
\begin{itemize}
    \item \textbf{Compression ratio calculation:} Compressed sizes were estimated using theoretical bit counts (e.g., 12‑bit codes for LZW and PSO). Actual bitstream sizes after entropy coding may differ slightly.
    \item \textbf{MSE vs. perceptual quality:} MSE and PSNR are imperfect proxies for visual quality, especially for colour quantisation. Part 2 will incorporate perceptual metrics (SSIM, MS‑SSIM) and user studies.
\end{itemize}

\subsection{Implications for Part 2}

The findings from Part 1 directly inform the continuation of this project in the next semester:

\begin{itemize}
    \item \textbf{Algorithm selection depends on use case:} LOCO‑I (JPEG‑LS) remains an excellent balance for lossless archival where speed is critical. Lossy+residual coding is superior for maximum lossless compression. For lossy colour quantisation, K‑Means provides fast, decent quality, while PSO can be reserved for offline quality‑critical tasks.
    \item \textbf{Hybrid methods are promising:} The success of lossy+residual and FBPNN suggests that combining complementary techniques (transform coding, prediction, fuzzy clustering, neural learning) can push compression beyond any single method.
    \item \textbf{Context matters:} The large performance gap between fixed‑predictor JPEG Lossless and context‑adaptive LOCO‑I underscores the importance of exploiting local image structure. Future codecs should incorporate adaptive prediction and context modelling as core components.
    \item \textbf{Optimisation can be traded for quality:} PSO improves palette design but at high computational cost. Part 2 will explore hybrid initialisation (K‑Means followed by a few PSO iterations) to balance quality and speed.
\end{itemize}

\subsection{Lessons Learned (Part 1)}

Several key insights were gained that will guide the next semester’s work:

\begin{itemize}
    \item \textbf{Simplicity is not always best:} Huffman coding is elegant but far from optimal for image compression due to its integer‑length limitation and lack of spatial awareness. Context‑based methods, though more complex, are well worth the implementation effort.
    \item \textbf{Numerical stability is critical:} Implementing arithmetic coding with integer arithmetic required careful handling of underflow and renormalisation. A seemingly small bug (e.g., an off‑by‑one error in interval scaling) can break losslessness. Thorough testing on small examples was invaluable.
    \item \textbf{Parameter tuning matters greatly:} PSO and Gath‑Geva performance was sensitive to the choice of \(w, c_1, c_2, m\). Default literature values worked reasonably, but systematic tuning would improve results.
    \item \textbf{Visualisation accelerates debugging:} Plotting residuals, histograms, and convergence curves helped identify problems (e.g., overflow in JPEG prediction) much faster than inspecting numerical logs.
    \item \textbf{Benchmarking is essential:} Including simple baselines (e.g., K‑Means for colour quantisation) from the start is good practice; it is impossible to claim a new method is better without direct comparison.
    \item \textbf{Fuzzy methods shine in ambiguous regions:} Fuzzy membership maps from Gath‑Geva revealed that pixels at colour boundaries (e.g., between skin and hair in Lenna) have high uncertainty. Weighting neural network updates by these memberships helped the FBPNN focus on clear examples, improving stability.
\end{itemize}

In summary, Part 1 of this project has demonstrated that no single compression algorithm dominates all scenarios. The optimal choice depends on the requirements: lossless vs. lossy, speed vs. quality, and the nature of the target image. Static methods are consistently outperformed by hybrid and adaptive approaches, and careful implementation – with attention to numerical stability and parameter tuning – remains essential to realise their full potential.

All results, lessons, and limitations identified here will serve as the foundation for **Part 2**, which will be carried out in the next semester. The continuation will focus on scaling to larger datasets, GPU acceleration, learned context models, perceptually guided optimisation, and end‑to‑end learned compression.